\section{Core Features}

% Structure this chapter by feature, not by one giant screenshot.
% Each feature can have multiple modules and each module should be explained
% using one container screenshot per card / table / form area.
% Use source code to identify what each module contains and what each
% container is called in the live UI.

\subsection{FEATURE\_NAME}

\begin{ugModule}{Module Name}
  \textbf{Purpose}

  Short explanation of what this module is for.

  \textbf{Who can access}

  State the roles or users that can open this module.

  \textbf{Key functions}
  \begin{itemize}
    \item Main action one.
    \item Main action two.
    \item Main action three.
  \end{itemize}

  \textbf{Step-by-step usage}
  \begin{ugSteps}
    \item Open the relevant menu item or container entry point.
    \ugScreenshot[0.55]{screenshots/feature-module-01.png}{The module entry container}
    \item Use the main control, button, tab, or filter inside that container.
    \ugScreenshot[0.35]{screenshots/feature-module-02.png}{The main control container}
    \item Fill in the required fields or review the list inside the module.
    \ugScreenshot[0.55]{screenshots/feature-module-03.png}{The form or table container}
    \item Save or submit your changes.
    \ugScreenshot[0.35]{screenshots/feature-module-04.png}{The save or submit container}
  \end{ugSteps}

  \textbf{Business rules}
  \begin{itemize}
    \item Rule one that users must follow.
    \item Rule two that affects validation or access.
    \item Rule three that affects the result.
  \end{itemize}

  \textbf{Expected result}

  Describe what should happen after the user completes the module.
\end{ugModule}

\begin{ugModule}{Another Module}
  Repeat the same structure for each module inside the feature.
\end{ugModule}

% Add a separate feature section for each major product area.
